\chapter{Rever Respostas e Verificar Resultados}
\label{ch:autocorrecao}

\ficcao O rascunho diz «faltou uma meia». O revisor troca a frase por «uma meia ficou por localizar». A segunda versão está mais cautelosa, mas ambos assumiram o mesmo número. Se o cliente falou de um par, precisamos de reler a mensagem.

\section{Pedir uma revisão com critério}
Podes pedir ao modelo que critique o seu rascunho e escreva uma nova versão. Define primeiro o que a revisão procura. Por exemplo, verifica a atribuição dos factos e corta duas frases repetidas. Numa tarefa de escrita, a segunda leitura dá outra oportunidade de encontrar um defeito.

\begin{lstlisting}
[FASE 1]
Responde à pergunta seguinte o melhor que conseguires.

[pergunta]

[FASE 2]
Agora lê a tua resposta anterior com espírito crítico.
Aponta:

possíveis erros factuais,

partes pouco claras,

oportunidades de melhoria.

[FASE 3]
Escreve uma nova versão da resposta incorporando
as correções que identificaste.
\end{lstlisting}

A revisão pode melhorar clareza. Huang e colaboradores encontraram casos de raciocínio em que a autocorreção sem feedback externo falhou ou piorou a resposta~\citep{huang2024}. Para verificar o número de peças, entrega ao revisor a passagem da reclamação e o cálculo de referência. O critério vem da tarefa. A apreciação do modelo é mais uma proposta a conferir.

\begin{promptlabbox}{Transformar o modelo em revisor de si próprio}
\textbf{Objetivo e contexto:} Rever uma resposta identificando primeiro defeitos concretos de clareza e suporte factual.

Escolhe uma resposta que o modelo te tenha dado recentemente e de que não tenhas gostado (texto confuso, muito genérico, etc.).

\medskip

\textbf{Passo 1} — Cola a resposta no chat e pede uma crítica detalhada.

\textbf{Passo 2} — Pede uma versão que corrija os defeitos identificados com apoio no texto de origem.

Compara a nova versão com a antiga. Assinala o que mudou na clareza e o que mudou nos factos, com a mesma grelha para ambas.

\textbf{Resultado esperado e validação:} Compara versões com a mesma grelha e confirma os factos numa fonte externa. A revisão pode melhorar a forma e piorar o conteúdo: regista ambos.

\textbf{Falhas possíveis:} Tratar uma crítica eloquente como prova, pedir apenas que fique ótima ou premiar mudanças sem ganho.
\end{promptlabbox}

\section{Protocolos iterativos: escrever \texorpdfstring{$\to$}{->} criticar \texorpdfstring{$\to$}{->} reescrever}

Podemos generalizar a ideia anterior em \emph{protocolos iterativos}. Em vez de uma única ronda de crítica, defines um processo:

\begin{enumerate}[leftmargin=*]
\item Gerar rascunho.
\item Criticar rascunho.
\item Reescrever com base na crítica.
\item (Opcional) Repetir 2–3 vezes, com critério de paragem.
\end{enumerate}

Para evitar loops intermináveis, convém pôr limites explícitos:

\begin{lstlisting}
Vamos seguir este processo:

Escreves um rascunho de resposta.

Escreves uma crítica ao rascunho, em 5 pontos.

Escreves uma nova versão melhorada.

Faz no máximo uma revisão. Se a fonte não sustentar
uma frase, assinala-a como pendente.

Começa pelo rascunho.
\end{lstlisting}

Podes até usar personas diferentes para cada fase:

\begin{itemize}[leftmargin=*]
\item Persona A: “autor”.
\item Persona B: “revisor mal-disposto”.
\item Persona C: “editor final”.
\end{itemize}

\begin{warningboxenv}{Mais iterações não significam melhor qualidade sem critério}
Se o critério de melhoria for vago (“melhora o texto”), o modelo pode ficar a girar em torno do mesmo estilo, mudando palavras sem ganhos reais. Usa críticas específicas (clareza, precisão, estrutura) para guiar a recursão.
\end{warningboxenv}

\section{Quando chamar Python, SQL ou outra ferramenta}

Há uma linha clara entre o que é cómodo fazer com LLMs e o que faz sentido delegar para ferramentas clássicas. Sinais de que está na hora de chamar Python, SQL ou outra coisa:

\begin{itemize}[leftmargin=*]
\item Precisão numérica é crítica (contabilidade, estatísticas, engenharia).
\item O problema envolve grandes volumes de dados estruturados.
\item Tens de repetir a mesma operação muitas vezes de forma idêntica.
\end{itemize}

Padrão de trabalho saudável:

\begin{enumerate}[leftmargin=*]
\item Usa o LLM para:
\begin{itemize}[leftmargin=*]
\item esclarecer o problema,
\item desenhar a estratégia de cálculo,
\item gerar código de esqueleto (em Python, SQL, etc.).
\end{itemize}
\item Usa a linguagem de programação / base de dados para:
\begin{itemize}[leftmargin=*]
\item executar os cálculos,
\item validar resultados,
\item automatizar pipelines.
\end{itemize}
\end{enumerate}

\begin{promptlabbox}{Fazer a ponte entre explicação e código}
\textbf{Objetivo e contexto:} Implementar um cálculo numérico num ambiente local, com inputs e resultados de referência definidos antes.

Escolhe um problema numérico simples mas chatinho (por exemplo, calcular prestações de um empréstimo com juros compostos).

\medskip

\textbf{Passo 1} — Pede ao LLM para explicar passo a passo como calcular.

\textbf{Passo 2} — Pede para gerar código Python que implemente esses passos.

\textbf{Passo 3} — Corre o código num ambiente real (fora do LLM) com diferentes valores e compara com resultados de uma calculadora independente.

Assim aproveitas o modelo como tutor/codificador, mas confias a execução à ferramenta certa.

\textbf{Resultado esperado e validação:} Executa o código, inclui valores de fronteira e compara com cálculo independente. Usa unidades e convenção de taxa explícitas; não testes com decisões financeiras reais.

\textbf{Falhas possíveis:} Confundir taxa anual com mensal, aceitar código não executado ou usar duas respostas do mesmo modelo como confirmação independente.
\end{promptlabbox}

\section{Prompting com ferramentas e \emph{function calling}}

Muitas aplicações ligam modelos a \emph{ferramentas} (APIs, bases de dados e calculadoras). Em function calling, o modelo propõe uma chamada estruturada; a aplicação valida argumentos, identidade e autorização antes de executar~\citep{functioncalling2026}. O serviço precisa de expor e autorizar a ferramenta. Com essa ligação, o modelo pode propor os passos seguintes.

\begin{itemize}[leftmargin=*]
\item interpretar o pedido,
\item decidir que ferramenta chamar,
\item construir os argumentos da chamada,
\item integrar o resultado da ferramenta na resposta final.
\end{itemize}

Do ponto de vista de prompting, isto muda o jogo:

\begin{itemize}[leftmargin=*]
\item Tens de explicar ao modelo \emph{quando} usar a ferramenta (não é óbvio).
\item Tens de definir claramente o que cada ferramenta faz e que argumentos aceita.
\end{itemize}

Um padrão típico (em pseudo-system prompt) é:

\begin{lstlisting}
Tens acesso às seguintes ferramentas:

get_exchange_rate(base, target)

Devolve a taxa de câmbio atual entre duas moedas.

consultar_estado_reclamacao(reclamacao_id)

Consulta o estado de uma reclamacao autorizada.
A identidade e as permissoes sao verificadas pelo servidor.

Usa uma ferramenta sempre que:

Precisares de dados atualizados.

Precisares de resultados numéricos exatos.

Quando a ferramenta devolver um resultado, explica
ao utilizador, em português simples, o que significa.
\end{lstlisting}

\begin{warningboxenv}{Valida no ponto em que a ação acontece}
Antes de consultar uma encomenda, verifica se o utilizador pode vê-la. Para alterar ou pagar, usa a autorização e os limites dessa operação. Um identificador bem formado é só um dos critérios. No próximo capítulo, vamos separar o formato dos dados da verificação do que afirmam.
\end{warningboxenv}

\section{Evitar loops infinitos e verificação inútil}

Protocolos iterativos e uso de ferramentas trazem um problema: se mal desenhados, podem cair em loops de:

\begin{itemize}[leftmargin=*]
\item reexplicar a mesma coisa,
\item pedir mais dados irrelevantes,
\item chamar ferramentas sem necessidade.
\end{itemize}

Regras simples para evitar isso:

\begin{itemize}[leftmargin=*]
\item Define \emph{critérios de paragem} claros nos prompts:
\begin{itemize}[leftmargin=*]
\item “No máximo 2 iterações de revisão.”
\item “Se não tiveres nova informação relevante, não faças nova tentativa.”
\end{itemize}
\item Distingue entre:
\begin{itemize}[leftmargin=*]
\item \emph{erro recuperável} (pode valer a pena tentar de novo),
\item \emph{erro estrutural} (falta de dados, pedido impossível).
\end{itemize}
\item Regista o nome da ferramenta, argumentos validados, resultado, duração e motivo operacional breve. O limite de passos e de tempo deve existir no executor, mesmo que o modelo não o cumpra.
\end{itemize}

\begin{keyideabox}{Mais passos não equivalem a mais inteligência}
É tentador montar pipelines com 10 fases “para garantir qualidade”. Sem critérios claros, isso só acrescenta latência, custo e pontos de falha. O objetivo é \emph{raciocínio suficiente}, não \emph{raciocínio infinito}.
\end{keyideabox}
